\chapter{Mô hình thị trường và hai tác tử}\label{chap:4}

Chương này hình thức hóa bài toán của Laboratory 02~\cite{lab02}: môi trường, trạng thái, hành động, phần
thưởng và hai chính sách. Trước khi chạy bất kỳ thí nghiệm nào, chương tính sẵn lợi nhuận kỳ vọng của mọi tổ hợp
và chính sách tối ưu, để các kết quả học được ở Chương~\ref{chap:6} có một mốc đối chiếu chính xác.

\section{Mô tả bài toán}

Hai cửa hàng thương mại điện tử cùng bán một loại sản phẩm. Ở mỗi bước thời gian $t$, mỗi cửa hàng $i\in\{1,2\}$
chọn đồng thời một hành động $a_i(t)\in\{0,1\}$: $0$ là giữ giá gốc, $1$ là giảm giá. Giảm giá kéo thêm khách
nhưng làm mỏng biên lợi nhuận của từng đơn, và nó còn lấy bớt khách của đối thủ. Cửa hàng~1 dùng tác tử học tăng
cường (Q-learning), cửa hàng~2 dùng tác tử theo luật cố định. Hình~\ref{fig:ch4-loop} tóm tắt một bước tương tác.

\hinhtikz{ch4_loop}{Vòng tương tác của hai tác tử trong một bước: (1) quan sát và chọn hành động đồng thời,
(2) thị trường tính xác suất chuyển đổi, (3) lấy mẫu doanh số và trả thưởng.}{ch4-loop}

\nhanxet{hai tác tử không trao đổi thông điệp; chúng ảnh hưởng nhau qua thị trường chung, đúng định nghĩa hệ đa tác
tử ở Chương~\ref{chap:1}. Tác tử luật cũng nhận phần thưởng nhưng không dùng nó để đổi hành vi.}

\section{Trạng thái và hành động}

Trạng thái cục bộ mà tác tử $i$ quan sát trước khi hành động là
\begin{equation}\label{eq:state}
  s_i(t) = \bigl(d_t,\; c_i(t)\bigr),\qquad d_t\in\{0,1\},\quad c_i(t) = a_{-i}(t-1)\in\{0,1\},
\end{equation}
trong đó $d_t$ là mức nhu cầu (0: thấp, 1: cao) và $c_i(t)$ là hành động của đối thủ ở bước trước. Mỗi tác tử có
$|\mathcal S| = 4$ trạng thái và $|\mathcal A|=2$ hành động, nên bảng $Q$ chỉ có $8$ ô. Trạng thái chung của cả
hệ là bộ ba $(d_t, a_1(t-1), a_2(t-1))$ với $8$ giá trị; tác tử chỉ thấy một phần của nó, cụ thể là không thấy
hành động trước của chính mình. Điểm này quan trọng ở Mục~\ref{sec:best}.

\section{Chuyển đổi khách hàng, doanh số và lợi nhuận}

Xác suất một khách tiềm năng mua hàng của cửa hàng $i$ là
\begin{equation}\label{eq:conv}
  p_i(t) = \operatorname{clip}\bigl(0{,}30 + 0{,}25\,d_t + 0{,}15\,a_i(t) - 0{,}10\,a_{-i}(t);\; 0{,}05;\; 0{,}95\bigr).
\end{equation}
Số khách tiềm năng là $N_t = 10$ khi $d_t = 0$ và $N_t = 20$ khi $d_t=1$; số đơn bán được
$S_i(t)\sim\mathrm{Binomial}(N_t, p_i(t))$. Với giá gốc $P=10$, giá vốn $C=5$ và mức giảm $\delta=0{,}20$, phần thưởng là
lợi nhuận
\begin{equation}\label{eq:reward}
  r_i(t) = S_i(t)\,\bigl[P\,(1-\delta\,a_i(t)) - C\bigr],
  \qquad \mathbb E[r_i\mid d,a_i,a_{-i}] = N_d\; p_i\; \bigl[P(1-\delta a_i)-C\bigr].
\end{equation}
Biên lợi nhuận mỗi đơn là $10-5=\SL{ch4_model/d20/margin_keep}$ khi giữ giá và $8-5=\SL{ch4_model/d20/margin_disc}$
khi giảm giá: giảm giá cắt $40\%$ biên lợi nhuận để đổi lấy $15$ điểm phần trăm xác suất chuyển đổi.

\maNguon{market__conversion_probability}{Xác suất chuyển đổi theo \eqref{eq:conv}.}

\section{Động lực môi trường}

Sau khi trả thưởng, môi trường sinh nhu cầu kế tiếp $d_{t+1}$ độc lập với hành động, đều trên $\{0,1\}$; hành động
hiện tại trở thành ``hành động trước'' của bước sau. Mỗi episode có $T=20$ bước; đầu episode nhu cầu được rút
ngẫu nhiên và hành động trước của cả hai đặt bằng $0$. Lớp \texttt{MarketConfig} còn có tham số
\texttt{demand\_persist} $=P(d_{t+1}=d_t)$: giá trị $0{,}5$ cho đúng mô hình đề bài, giá trị $0{,}8$ cho biến thể
``nhu cầu bền'' dùng ở câu~9. Mã của lớp môi trường được trình bày ở Chương~\ref{chap:5}.

\section{Tác tử Q-learning}

Tác tử RL giữ bảng $Q:\mathcal S\times\mathcal A\to\mathbb R$. Sau mỗi bước nó cập nhật cặp vừa dùng~\cite{ch2:watkins}:
\begin{equation}\label{eq:qupdate}
  Q(s_t,a_t) \leftarrow Q(s_t,a_t) + \alpha\Bigl[\underbrace{r_t + \gamma\max_{a'}Q(s_{t+1},a')}_{\text{mục tiêu TD}} - Q(s_t,a_t)\Bigr],
\end{equation}
với $\alpha = 0{,}10$, $\gamma=0{,}90$; ở bước cuối episode mục tiêu chỉ còn $r_t$. Hành động được chọn theo
$\varepsilon$-greedy: ngẫu nhiên với xác suất $\varepsilon$, ngược lại chọn $\arg\max_a Q(s,a)$ (hòa thì chọn ngẫu
nhiên), với lịch giảm dần
\begin{equation}\label{eq:eps}
  \varepsilon_e = \max\bigl(\varepsilon_{\min},\; \varepsilon_0\,\rho^{\,e}\bigr) = \max\bigl(0{,}02;\; 0{,}20\cdot 0{,}995^{\,e}\bigr)
\end{equation}
ở episode $e$. Sau khoảng $460$ episode $\varepsilon$ chạm sàn $0{,}02$ và giữ nguyên tới cuối.

\section{Tác tử theo luật}

Tác tử luật dùng chính sách if--else viết sẵn (Hình~\ref{fig:ch4-rule}):
\begin{equation}\label{eq:rule}
  \pi_{\text{luật}}(d,c) = \begin{cases} 0, & d=0,\\ 1, & d=1 \text{ và } c=1,\\ 0, & d=1 \text{ và } c=0.\end{cases}
\end{equation}
Nhu cầu thấp thì giữ giá; nhu cầu cao thì chỉ giảm giá để ``trả đũa'' khi đối thủ vừa giảm.

\hinhtikz{ch4_rule}{Cây quyết định của tác tử theo luật.}{ch4-rule}

\section{Phân tích lợi nhuận kỳ vọng một bước}\label{sec:payoff}

Vì mô hình nhỏ, có thể tính chính xác \eqref{eq:reward} cho mọi tổ hợp nhu cầu $d$, hành động của mình và hành động
\emph{hiện tại} $b$ của đối thủ (Bảng~\ref{tab:ch4-reward}).

\bang{ch4_reward}{Lợi nhuận kỳ vọng một bước của cửa hàng $i$ theo $(d,b)$, với $\delta=20\%$ và $\delta=5\%$.}{ch4-reward}

\nhanxet{với $\delta=20\%$, giữ giá tốt hơn ở ba trên bốn ô. Chỉ ở ô $(0,1)$, khi nhu cầu thấp và đối thủ đang giảm
giá, giảm giá nhỉnh hơn ($10{,}5$ so với $10$). Với $\delta=5\%$ biên lợi nhuận khi giảm còn $4{,}5$, và giảm giá thắng
ở mọi ô có nhu cầu thấp.}

Hình~\ref{fig:ch4-delta} quét $\delta$ từ $0$ đến $40\%$: đường liền (giảm giá) cắt đường đứt (giữ giá) tại các ngưỡng
khác nhau theo từng trạng thái. Đây là lý do Thí nghiệm F (Chương~\ref{chap:6}) đổi $\delta$ để xem kết luận ``RL hay luật
tốt hơn'' có đổi chiều hay không.

\hinhpdf{ch4_reward_delta}{1}{Lợi nhuận kỳ vọng một bước theo mức giảm giá $\delta$; đường chấm đứng là $\delta=20\%$.}{ch4-delta}

\section{Chính sách tối ưu khi đối thủ là tác tử luật}\label{sec:best}

Bảng một bước chưa đủ, vì hành động hôm nay đổi trạng thái của đối thủ ngày mai. Nếu RL giảm giá ở bước $t$ và
nhu cầu ở $t+1$ cao, tác tử luật sẽ giảm giá ở $t+1$ và lợi nhuận của RL tụt từ $55$ xuống $45$. Để có mốc chính xác,
hàm \texttt{exact\_value} lan truyền phân phối của trạng thái chung qua $20$ bước và cộng phần thưởng kỳ vọng, không
lấy mẫu; \texttt{best\_stationary} liệt kê toàn bộ $2^4=16$ chính sách tất định trên $(d,c)$; \texttt{optimal\_value}
giải quy hoạch động ngược trên trạng thái chung và bước $t$ (cận trên thật).

\begin{phathien}{Kết quả phân tích (output của \texttt{scripts/summarize.py})}
\begin{itemize}
  \item Hai tác tử luật đối đầu: không ai giảm giá (hành động trước ban đầu bằng $0$), mỗi bên đạt
    $\SL{ch4_model/d20/rule_rule}$ mỗi bước.
  \item Tối ưu trên 16 chính sách dừng: \texttt{\SL{ch4_model/d20/best_policy}} (không bao giờ giảm giá), giá trị
    $\SL{ch4_model/d20/best_value}$; tối ưu quy hoạch động: $\SL{ch4_model/d20/optimal}$. Tác tử luật vì vậy đã tối
    ưu ở cấu hình tham chiếu; RL giỏi nhất cũng chỉ hòa.
  \item Với $\delta=5\%$: tối ưu là \texttt{\SL{ch4_model/d05/best_policy}} (luôn giảm giá), giá trị
    $\SL{ch4_model/d05/optimal}$, cao hơn luật ($\SL{ch4_model/d05/rule_rule}$). Ở đây RL có dư địa để thắng.
\end{itemize}
\end{phathien}

\begin{luuy}{Vì sao đây là bài toán đa tác tử khó}
Hành động của đối thủ có mặt trong xác suất chuyển đổi, trong trạng thái và gián tiếp trong phần thưởng. Trạng thái
$(d,c)$ thiếu $a_i(t-1)$, thứ quyết định hành động \emph{hiện tại} của tác tử luật, nên quá trình RL thấy không còn
Markov hoàn toàn; khi đối thủ cũng học, môi trường còn trở thành phi dừng~\cite{busoniu2008}.
\end{luuy}
